\documentclass[a4paper,11pt,fleqn]{article}%fleqn pour aligner les equations à gauche

\usepackage{subcaption}
\usepackage{graphicx}
\usepackage{wrapfig}

\usepackage{romannum}

\usepackage[utf8]{inputenc}  
\usepackage[T1]{fontenc}
\usepackage{lscape}

\usepackage{tgheros,textcomp}%Police Arial
\renewcommand{\familydefault}{\sfdefault}

%\setlength{\mathindent}{0pt}%indentation dans math = 0

\usepackage[top = 1cm, bottom = 2cm, left = 1cm, right = 1cm]{geometry}
\usepackage{titlesec}
%\setcounter{secnumdepth}{4}
\usepackage{xcolor}
\usepackage{amssymb} %double flèche de réaction
\usepackage{mathrsfs}%farraday
\usepackage{xtab}%tableau sur plusieurs pages
\usepackage{esvect}%grands vecteurs

\usepackage[version=4]{mhchem} %equations chimiques

\usepackage{array}
\usepackage{chemfig}
\usepackage{multirow}
\usepackage{sectsty}

%\usepackage{caption}

%\usepackage{adjustbox}
\usepackage[labelformat=simple]{subfig}


\usepackage{pgfmath}
\usepackage{amsmath}
\usepackage{amssymb}
\usepackage[cal=boondoxo]{mathalfa} 
\usepackage{enumitem}
\usepackage{lastpage}
\usepackage{tcolorbox}
\usepackage{siunitx}
\usepackage[european, straightvoltages, RPvoltages]{circuitikz}
\usepackage{tikz}
\usetikzlibrary{shapes.geometric}
\usetikzlibrary{calc}
\usetikzlibrary{automata,positioning}
\usetikzlibrary{shadows}
\usepackage{multicol}

\usepackage{tabularx,colortbl}%colorer fond cellules

\sisetup{
	detect-all,
	output-decimal-marker={,},
	group-minimum-digits = 3,
	group-separator={~},
	number-unit-separator={~},
	inter-unit-product={~}
} %setup pour SIunitx

\usepackage{listings}%pour insérer le code python
\definecolor{codegreen}{rgb}{0,0.6,0}
\definecolor{codegray}{rgb}{0.5,0.5,0.5}
\definecolor{codepurple}{rgb}{0.58,0,0.82}
\definecolor{backcolour}{rgb}{0.95,0.95,0.92}
\definecolor{orangecirc} {rgb}{231, 149, 84}
\definecolor{greencirc} {rgb}{47, 166, 55}



\titleformat{\section}
{
\vspace{.3ex}%
\normalfont \bfseries \large}
{\thesection.}{.5em}{\bfseries}

\renewcommand{\thesection}{\arabic{section}}
\titlespacing{\section}{0pt}{1em}{.5em}
\renewcommand{\thesubsection}{\hspace{1cm}\alph{subsection}.}
\renewcommand\thesubfigure{\textbf{Document \arabic{subfigure} :}}

\title{\vspace{-1cm}\large 
	\begin{tabular}{|m{5cm}|>{\centering}m{8cm}|>{\centering\arraybackslash}m{4cm}|}
		\hline
		Nom : & \textbf{\textsc{DS N°5}} & Note : \\
		Prénom : & & \\\hline
\end{tabular}}
\date{}

\newpagestyle{mystyle}{%
	\setfoot{ROSALIE}{\thepage\,$|$\,\pageref{LastPage}}{}
}
\renewpagestyle{plain}{%
	\setfoot{ROSALIE}{\thepage\,$|$\,\pageref{LastPage}}{}
}

\pagestyle{mystyle}

\newenvironment{itemize*}%
{\begin{itemize}[label=$-$,labelindent=0pt, leftmargin=*,topsep=0pt]%
		\setlength{\itemsep}{0pt}%
		\setlength{\parskip}{0pt}}%
	{\end{itemize}}

\newenvironment{itemize**}%
{\begin{itemize}[label=,labelindent=0pt, leftmargin=*,topsep=0pt]%
		\setlength{\itemsep}{0pt}%
		\setlength{\parskip}{0pt}}%
	{\end{itemize}}

\renewcommand{\arraystretch}{1.5}
\renewcommand{\figurename}{Figure}
  
\begin{document}
\maketitle
\vspace{-2.5cm}


\section{Une potion radioactive}
De tous temps, certaines substances sont considérées comme des remèdes contre tous les maux, des panacées.
\begin{minipage}{.7\textwidth}
    Au début du \Romannum{20} siècle, le Radithor, sorte de "potion magique" était censé soigner plus d'une centaine de maladies.\\
    Un cancérologue américain a trouvé chez un antiquaire, plusieurs bouteilles de Radithor. Bien que vidées depuis 10 ans de leur contenu, les bouteilles se sont avérées être encore dangereusement radioactives. Chacune avait vraisemblablement contenu environ un microcurie de radium 226 et de radium 228\\
    1 microcurie correspond à \SI{3,7E4}{Bq}.
\end{minipage}
\begin{minipage}{.3\textwidth}
    \centering
        \includegraphics[width=.5\linewidth]{Ressources/TSPE_DS5_Fig1.jpg}
\end{minipage}

\begin{tabular}{|>{\centering}m{3cm}|c|c|c|c|c|}
     \hline \textbf{Noyau}&Radium 226&Radium 228&Actinium 228&Radon 222&Hélium 4\\
    \hline \textbf{Symbole}&\ce{^226_88Ra}&\ce{^228_88Ra}&\ce{^228_89Ac}&\ce{^222_86Rn}&\ce{^4_2He}\\
    \hline
\end{tabular}\\

\begin{tabular}{|>{\centering}m{3cm}|c|c|c|}
     \hline \textbf{Noyau}&Radium 226&Radon 222&Hélium 4\\
    \hline \textbf{Masse en $u$}&\num{225,9970}&\num{221,9703}&\num{4,0015}\\
    \hline
\end{tabular}
\begin{itemize*}
    \item \textbf{Unité de masse atomique} : $\SI{1}{u}=\SI{1,66054E-27}{kg}$
    \item \textbf{Célérité de la lumière dans le vide} : $c=\SI{3,00E8}{\meter\per\second}$
    \item \textbf{Constante d'Avogadro} : $N_A=\SI{6,02E23}{\per\mole}$
    \item \textbf{Masse molaire du \ce{^226_88Ra}} : $M=\SI{226}{\gram\per\mole}$
\end{itemize*}


\subsection{Le radium 226 et le médothorium}
Sur l'étiquette du flacon de Radithor est mentionnée la présence de mésothorium, ancienne dénomination du radium 228. Cette "eau certifiée radioactive" contenait également du radium 226.
\begin{enumerate}
    \item Les noyaux de radium 228 et de radium 226 sont des isotopes. Expliquer.
    \item Le radium 228 se désintègre pour donner l'isotope 228 de l'actinium \ce{Ac} et une particule notée \ce{X}. $$\ce{^228_88Ra -> ^228_89Ac + ^A_Z X}$$
    Compléter l'équation de désintégration en citant les lois utilisées puis identifier \ce{X}. De quel type de radioactivité s'agit-il ?
\end{enumerate}
Dans la suite de l’exercice, on néglige la présence de radium 228 dans le Radithor.\\ 
On suppose que l’activité radioactive du flacon est uniquement due à la présence de l’isotope 226 du radium. Celui-ci se désintègre spontanément selon l’équation suivante :
$$\ce{^226_88Ra -> ^222_86Ac + ^4_2He}$$

\subsection{Constante radioactive du radium 226}
 L’activité A(t) d’un échantillon de noyaux de radium 226 suit la loi de décroissance exponentielle $A(t)=A_0\cdot e^{-\lambda \cdot t}$, avec $A_0$ l'activité à $t=\SI{0}{s}$.
 \begin{enumerate}[resume]
     \item Rappeler la définition de la demi-vie $t_{1/2}$ d'un échantillon radioactif.
     \item Vérifier sur la courbe donnant l'évolution de l'activité de l'échantillon en fonction du temps représentée dans le \textbf{Figure~\ref{doc1}}, que la demi-vie du radium 226 est égale à \SI{1,60E3}{ans}.
     \item Établir la relation entre la demi-vie et la constante radioactive $\lambda$ puis calculer la valeur de $\lambda$ en \unit{\per\second}.
 \end{enumerate}
 \subsection{Masse de radium 226}
 \begin{enumerate}[resume]
     \item Donner la relation liant l'activité $A(t)$ d'un échantillon radioactif au nombre $N(t)$ de noyaux radioactifs présents.
     \item Calculer $N_0$, le nombre de noyaux de radium 226 initialement présents dans le flacon de Radithor.
     \item Vérifier que le flacon contenait alors une masse $m_0 = \SI{1,0}{\micro\gram}$ de radium 226.
 \end{enumerate}

 \begin{figure}[h!]
    \caption{Activité d’un échantillon de radium 226 en fonction du temps }
    \centering
    \includegraphics[width=.9\linewidth]{Ressources/TSPE_DS5b_Fig1.png}
    \label{doc1}
 \end{figure}
\end{document}

